\section{因果识别的尝试与失败}\label{sec:identification}

\subsection{设定与样本}\label{sec:id-sample}

以城市 $\times$ 春节年份为观测单位，处理变量为“是否全面禁放”（\texttt{ban}）。因中国禁放政策交错实施、且 2023 年后部分城市由禁改限，处理\textbf{可逆转}，故主估计量采用 de Chaisemartin \& D'Haultf\oe uille（2020）允许处理逆转的估计量；Callaway \& Sant'Anna（2021）要求处理为吸收态，仅作附录对照。多向固定效应用交替投影（FWL）实现。

样本为 \textbf{306 城 $\times$ 6 年 $=1{,}836$ 行}（2021—2026）。窗口起点取 2021，是因为山西各市的首条政策记录始于该年；更早的窗口会把“首次进入记录”误当成一次 $0\rightarrow1$ 的切换。

三个批次，方向不同：

\begin{table}[H]
\centering
\caption{三个政策批次}
\label{tab:batches}
\small
\begin{tabular}{L{1.3cm}L{7.0cm}L{2.2cm}L{2.4cm}R{1.6cm}}
\toprule
\textbf{批次} & \textbf{城市} & \textbf{方向} & \textbf{同省对照} & \textbf{处理单位} \\
\midrule
2023 & 信阳、南阳、平顶山、洛阳、郑州 & 退出（禁$\rightarrow$限） & 河南其余 12 城 & 5 \\
2024 & 东营、滨州 & \textbf{进入}（限$\rightarrow$禁） & 山东其余城市 & 2 \\
2026 & 临汾、吕梁、大同、太原、晋中、晋城、朔州、运城、长治、阳泉 & 退出（省级统一） & \textbf{0} & \textbf{1} \\
\bottomrule
\end{tabular}
\end{table}

\noindent 2024 年批次是样本期内\textbf{唯一的收紧方向}变异：东营依据 2024-01-02 施行的政府令第 208 号转为全年禁放（第~\ref{sec:institution}~节）；滨州同日施行，但该条证据档为 \texttt{verified=N}、原文未取得。

\paragraph{处理变量的覆盖是本节最根本的限制。}

政策库覆盖 34 城，其余 \textbf{272 城没有任何政策记录}。在 3,672 个城市-年中，\texttt{policy\_type} 非空者仅 \textbf{243 个（6.6\%）}；限定到估计窗口 2021—2026，1,836 行中只有 \textbf{201 行（10.9\%）}。代码对空值的处理是“按 0（未禁放）”，但这是一条\textbf{记录口径，不是事实}：一个没有政策记录的城市完全可能一直在禁放。

因此本文所谓的“全国对照”，实际是“没有政策记录的城市”；而 2019—2022 年的大面积禁放使这个对照本身被污染。测量误差的方向是衰减——它不能救回设计，但也不能被忽略。\textbf{这一点比第~\ref{sec:why-fail}~节讨论的同文下达问题更基本。}

\subsection{诊断}

\paragraph{诊断零：安慰剂结果变量不通过。}

用节前基线（除夕前 7 日至前 2 日）作结果变量。燃放管制不应改变年前的浓度。

\begin{table}[H]
\centering
\caption{安慰剂结果变量：节前基线}
\label{tab:placebo-outcome}
\small
\begin{tabular}{L{4.4cm}R{2.0cm}R{2.4cm}R{2.4cm}}
\toprule
\textbf{设定} & $\beta$ & \textbf{SE} & \textbf{有效聚类数} \\
\midrule
城市 $+$ 年份 & $\mathbf{+10.62}$ & 2.54 & 306 \\
城市 $+$ 省份 $\times$ 年份 & $\mathbf{-4.48}$ & 1.59 & 33 \\
\bottomrule
\end{tabular}
\end{table}

两个设定下都显著，且方向相反。\textbf{“是否禁放”能预测节前的基线浓度}——这直接说明处理变量与城市层面的长期趋势相关，而不是与燃放相关。

\paragraph{诊断一：安慰剂“日”扫描。}

把结果变量换成窗口内的每一个相对日，跑同一个双重差分。若估计量锁定在燃放上，应当只有初一显著。

\begin{table}[H]
\centering
\caption{安慰剂“日”扫描（选取的若干日）}
\label{tab:placebo-days}
\small
\begin{tabular}{lR{1.3cm}R{1.3cm}R{1.6cm}R{1.3cm}R{1.3cm}}
\toprule
\multicolumn{6}{l}{\textbf{河南（同省对照，17 城，2021—2026）}} \\
\texttt{rel\_day} & $-6$ & $-4$ & $\mathbf{+1}$ & $+3$ & $+6$ \\
$\beta$ & $+18.7$ & $-15.0$ & $\mathbf{-10.8}$ & $+35.3$ & $-14.2$ \\
$t$ & $+2.38$ & $-2.60$ & $\mathbf{-0.56}$ & $+3.21$ & $-3.15$ \\
\midrule
\multicolumn{6}{l}{\textbf{山西 2026 放开（全国对照）}} \\
\texttt{rel\_day} & $-4$ & $-2$ & $\mathbf{+1}$ & $+5$ & $+6$ \\
$\beta$ & $-28.0$ & $-20.8$ & $\mathbf{+37.7}$ & $+24.9$ & $+54.6$ \\
$t$ & $-8.03$ & $-6.76$ & $\mathbf{+3.02}$ & $+2.81$ & $+4.49$ \\
\bottomrule
\end{tabular}
\end{table}

河南：\textbf{初一不显著（$t=-0.56$），而与燃放无关的 $-6$、$-4$、$+3$、$+6$ 都显著。}

山西：初一显著（$+37.7$，$t=3.02$），方向与“放开后上升”一致；\textbf{但除夕前一周的每一天都显著为负}（$-7$ 到 $-1$ 的 $t$ 值在 $-2.8$ 到 $-8.0$ 之间）。这不是一个干净的对照：处理组在政策变动之前就在相对下行。

\paragraph{诊断二：事件路径的前置期不平稳。}

三个批次的前置期（原始差，处理组变化减对照组变化）：

\begin{table}[H]
\centering
\caption{各批次事件路径的前置期}
\label{tab:preperiods}
\small
\begin{tabular}{L{2.0cm}L{8.8cm}R{2.4cm}}
\toprule
\textbf{批次} & \textbf{前置期} & \textbf{切换当期} \\
\midrule
2023 河南 & $-1$：$-46.61$ & $+75.30$ \\
2024 山东 & $-2$：$-56.87$；$-1$：$+141.97$ & $-81.48$ \\
2026 山西 & $-4$：$-45.83$；$-3$：$+32.42$；$-2$：$+15.24$；$-1$：$-22.33$ & $+44.61$ \\
\bottomrule
\end{tabular}
\end{table}

三个批次的前置期都出现与“效应”同量级甚至更大的波动。\textbf{平行趋势不成立。}

\paragraph{诊断三：量级先验不允许 30—50 的效应。}

中位城市的整个春节尖峰只有 28.1 \ugm（平衡面板口径）。声称禁放削减 30—50，等于消灭中位城市尖峰的 107\%—178\%。

\subsection{固定效应设定与窗口敏感性}

\begin{table}[H]
\centering
\caption{两个固定效应设定}
\label{tab:fe}
\small
\begin{tabular}{L{4.4cm}R{2.0cm}R{2.2cm}R{2.2cm}R{2.4cm}}
\toprule
\textbf{设定} & $\beta$ & \textbf{SE（CR1）} & \textbf{SE（含校正）} & \textbf{有效聚类数} \\
\midrule
城市 $+$ 年份 & $-31.74$ & 10.60 & 11.6—12.4 & 306 \\
城市 $+$ 省份 $\times$ 年份 & $-5.34$ & 14.41 & 14.8—15.2 & \textbf{33} \\
\bottomrule
\end{tabular}
\end{table}

省份 $\times$ 年份设定下只有 \textbf{33 个城市}的残差化处理变量非零，其余 273 城的贡献为零。窗口敏感性：$\beta$ 在 \texttt{min\_year} 取 2019、2021、2022 时，城市 $+$ 年份设定为 $-40.62$、$-31.74$、$-40.48$；省份 $\times$ 年份设定为 $-8.09$、$-5.34$、$-7.62$。\textbf{两个设定的符号一致，但量级随窗口变动约 40\%。}

\begin{quote}
\textbf{一处更正（2026-10-20）。} 论文的上一版报告 $\beta=-42.60$（SE 9.52）与 $+10.80$（SE 18.77，有效聚类 17），并据此写出“\textbf{符号在两个设定之间翻转}”这一诊断。那些数字是在政策库尚未纳入 2024 年山东批次时跑的。纳入东营、滨州两条限$\rightarrow$禁记录后，\textbf{符号翻转不再出现}，该诊断已撤回。

撤回它并不削弱本文的结论，因为结论依赖的是本节的诊断零至诊断三与处理变量覆盖，而不是符号翻转。但“\textbf{小标准误不等于可信}”这条方法论教训仍然成立，只是证据换成了：安慰剂结果变量显著（$+10.62$，$t=4.2$）、事件路径前置期剧烈波动、以及 93\% 的城市-年没有处理记录。
\end{quote}

\subsection{随机化推断}

在河南 17 城、5 个处理单位的设计下做精确随机化推断（$C(17,5)=6{,}188$ 种分配）。观测统计量 $-26.20$，双侧 $p=0.1836$，位于零分布的 $-1.35$ 个标准差处，零分布标准差 19.39。\textbf{不能拒绝零假设。}

\paragraph{校准性检查。}

一个常见的错误论证是：“假切换期的 $p=0.85$，说明检验是校准的。”\textbf{这个理由不成立}——校准的检验在零假设下 $p$ 值服从 $U(0,1)$，观察到 0.85 与观察到 0.15 同样可能。正确的检查是对河南 8 次\textbf{无政策变动}的转移逐次跑同一检验，得 $p=0.043$、$0.173$、$0.447$、$0.516$、$0.587$、$0.850$、$0.929$、$0.999$，\textbf{KS 距离 0.225}（$n=8$ 的 5\% 临界值约 0.457，不拒绝）。校准性未被推翻，但依据是这一条。

\paragraph{边界。} 本检验回答的是“哪几个城市改了政策”相对于省内随机分配是否罕见，\textbf{它不是因果零假设下的检验}——政策内生正是主要担忧。它真正有价值的产出是噪声尺度，而不是那个 $p$ 值。

\subsection{为什么失败，以及这不等于什么}\label{sec:why-fail}

失败有两层原因。

\textbf{第一层是测量：93\% 的城市-年没有处理记录。} 这不是编码问题，是政策库的覆盖问题。“全国对照”实为“无政策记录的城市”。

\textbf{第二层是共变：禁令几乎总是与同期的一揽子工业减排措施同文下达。} 东营的链路是明证：禁令写在一份以生态环境部门为主、与量化工业减排指标同文下达的方案里。实测显示 \NOtwo{} 与 \SOtwo{} 在年前 6 天就已显著下降，初一反而不显著——指向工业与采暖，而不是烟花。

（滨州 2024 年通常被认为同理，但其方案原文未取得、证据档为 \texttt{verified=N}，\textbf{本文不把它当作已核实的案例}。）

因此本文必须把两句话分开写：

\begin{itemize}
  \item \textbf{正确}：“这个设计无法把禁令从同期的工业减排中分离出来。”
  \item \textbf{错误}：“禁放没有效果。”
\end{itemize}

后者是对数据的误读，也是本文拒绝作出的推断。既有文献（陆卜梵、刘蓓蓓，2022；Chen et al., 2022）在更有利的设计下确实报告了显著效应，本文没有推翻它们，本文只是没能复现其识别。

\subsection{关于三轮独立复核}\label{sec:replication}

2026-10-15，本文安排了三轮彼此独立的复核（隔离约束与结果见附录~\ref{app:replication}）。复核在\textbf{当时的面板}上逐位复现了六个数字，并排除了“管线有 bug 制造了零结果”这一替代解释。

\textbf{但那些数字属于旧面板。} 政策库纳入 2024 年山东批次后，主估计变为 $-31.74$ 与 $-5.34$，复核过的具体数值不再适用于本版。复核所验证的三件事仍然成立：交替投影的实现正确、\texttt{rel\_day} 映射正确、政策库的展开正确。\textbf{本版的所有估计均由当前面板重新跑出，复现命令见附录~\ref{app:identification}。}
